\subsection{Good Friday}
\fancyhead[RO,LE]{\textit{Good Friday}}
\fancyhead[RE,LO]{}
\begin{inhead}
    {First Class Double}
\end{inhead}

\begin{rubric}
	In Choir, None being ended, the Priest and Ministers in black vestments, without lights and incense, proceed to the Altar: and pray awhile prostrate before it. Meanwhile the Acolytes spread one cloth only upon the Altar. The Priest with the Ministers, having finished their prayer, go up to the Altar, and he kisses it in the middle: then the Reader proceeds to read the Prophecy in the place where the Epistle is read, and begins it without any announcement: the Priest meanwhile reads the same in a low voice at the Epistle corner.
\end{rubric}

\readingcitation{Lesson}{Hosea 6:1}
\lett{T}{hus saith the Lord:} In their affliction they will seek me early: Come, and let us return unto the \divineName{Lord}: for he hath torn, and he will heal us; he hath smitten, and he will bind us up. After two days will he revive us: in the third day he will raise us up, and we shall live in his sight. Then shall we know, if we follow on to know the \divineName{Lord}: his going forth is prepared as the morning; and he shall come unto us as the rain, as the latter and former rain unto the earth. O Ephraim, what shall I do unto thee? O Judah, what shall I do unto thee? for your goodness is as a morning cloud, and as the early dew it goeth away. Therefore have I hewed them by the prophets; I have slain them by the words of my mouth: and thy judgments are as the light that goeth forth. For I desired mercy, and not sacrifice; and the knowledge of God more than burnt offerings.

\begin{rubric}
	\blackrubric{Thanks be to God} is not answered here, nor after the following Lesson.
\end{rubric}

\tract{O Lord, I have heard thy speech, and was afraid: I have considered thy works, and was confounded. ℣. O Lord, revive thy work in the midst of the years: in the midst of the years make it known. ℣. In the time of confusion of my soul: in wrath, remember mercy. ℣. God came from Teman, and his Holy One from the thick woods of the mountains. ℣. His glory covered the heavens: and the earth was full of his praise.}

\begin{rubric}
	The Tract being finished, the Priest at the Epistle corner says: \blackrubric{Let us pray.} The Deacon: \blackrubric{Let us bow the knee,} and the Subdeacon: \blackrubric{Arise.}
\end{rubric}

%MANUAL ADJUSTMENT:
\clearpage
\collect
\lett{O}{God,} from whom Judas received the punishment of his guilt, and the thief the reward of his confession, grant unto us the effects of thy propitiation: that as in his passion Jesus Christ, our Lord, gave unto both the divers rewards of their merits; so he may take away the transgressions of our old nature, and bestow upon us the grace of his resurrection. Who liveth.

\readingcitation{Lesson}{Hebrews 10:1}

\begin{rubric}
	The Subdeacon in the tone of the Epistle, likewise without announcement, sings the following Lesson:
\end{rubric}

\lett{B}{rethren:} The law having a shadow of good things to come, and not the very image of the things, can never with those sacrifices which they offered year by year continually make the comers thereunto perfect. For then would they not have ceased to be offered? because that the worshippers once purged should have had no more conscience of sins. But in those sacrifices there is a remembrance again made of sins every year. For it is not possible that the blood of bulls and of goats should take away sins. Wherefore when he cometh into the world, he saith, Sacrifice and offering thou wouldest not, but a body hast thou prepared me: in burnt offerings and sacrifices for sin thou hast had no pleasure. Then said I, Lo, I come (in the volume of the book it is written of me,) to do thy will, O God. Above when he said, Sacrifice and offering and burnt offerings and offering for sin thou wouldest not, neither hadst pleasure therein; which are offered by the law; then said he, Lo, I come to do thy will, O God. He taketh away the first, that he may establish the second. By the which will we are sanctified through the offering of the body of Jesus Christ once for all. And every priest standeth daily ministering and offering oftentimes the same sacrifices, which can never take away sins: but this man, after he had offered one sacrifice for sins for ever, sat down on the right hand of God; from hence-forth expecting till his enemies be made his footstool. For by one offering he hath perfected for ever them that are sanctified. Whereof the Holy Ghost also is a witness to us: for after that he had said before, This is the covenant that I will make with them after those days, saith the Lord, I will put my laws into their hearts, and in their minds will I write them; and their sins and iniquities will I remember no more. Now where remission of these is, there is no more offering for sin. Having therefore, brethren, boldness to enter into the holiest by the blood of Jesus, by a new and living way, which he hath consecrated for us, through the veil, that is to say, his flesh; and having an high priest over the house of God; let us draw near with a true heart in full assurance of faith, having our hearts sprinkled from an evil conscience, and our bodies washed with pure water. Let us hold fast the profession of our faith without wavering; (for he is faithful that promised;) and let us consider one another to provoke unto love and to good works: not forsaking the assembling of ourselves together, as the manner of some is; but exhorting one another: and so much the more, as ye see the day approaching.

\tract{Deliver me, O Lord, from the evil man: and preserve me from the wicked man. ℣. Who imagine mischief in their hearts: and stir up strife all the day long. ℣. They have sharpened their tongues like a serpent: adders' poison is under their lips. ℣. Keep me, O Lord, from the hands of the ungodly: and preserve me from the wicked men. ℣. Who are purposed to overthrow my goings: the proud have laid a snare for me. ℣. And spread a net abroad with cords: yea, and set traps in my way. ℣. I said unto the Lord: Thou art my God: hear the voice of my prayers, O Lord. ℣. O Lord God, thou strength of my health: thou hast covered my head in the day of battle. ℣. Let not the ungodly have his desire, O Lord: let not his mischievous imagination prosper, lest they be too proud. ℣. Let the mischief of their own lips fall upon the head of them: that compass me about. ℣. The righteous also shall give thanks unto thy name: and the just shall continue in thy sight.}

\begin{rubric}
	The Tract ended, the Gospel is read upon a bare pulpit: the Celebrant meanwhile reads it in a low voice at the Epistle corner.
\end{rubric}

\readingcitation{Gospel}{John 19:1}
\lett{A}{t that time:} Pilate took Jesus, and scourged him. And the soldiers platted a crown of thorns, and put it on his head, and they put on him a purple robe, and said, Hail, King of the Jews! and they smote him with their hands. Pilate therefore went forth again, and saith unto them, Behold, I bring him forth to you, that ye may know that I find no fault in him. Then came Jesus forth, wearing the crown of thorns, and the purple robe. And Pilate saith unto them, Behold the man! When the chief priests therefore and officers saw him, they cried out, saying, Crucify him, crucify him. Pilate saith unto them, Take ye him, and crucify him: for I find no fault in him. The Jews answered him, We have a law, and by our law he ought to die, because he made himself the Son of God. When Pilate therefore heard that saying, he was the more afraid; and went again into the judgment hall, and saith unto Jesus, Whence art thou? But Jesus gave him no answer. Then saith Pilate unto him, Speakest thou not unto me? knowest thou not that I have power to crucify thee, and have power to release thee? Jesus answered, Thou couldest have no power at all against me, except it were given thee from above: therefore he that delivered me unto thee hath the greater sin. And from thenceforth Pilate sought to release him: but the Jews cried out, saying, If thou let this man go, thou art not C{\ae}sar's friend: whosoever maketh himself a king, speaketh against C{\ae}sar. When Pilate therefore heard that saying, he brought Jesus forth, and sat down in the judgment seat in a place that is called the Pavement, but in the Hebrew, Gabbatha. And it was the preparation of the passover, and about the sixth hour: and he saith unto the Jews, Behold your King! But they cried out, Away with him, away with him, crucify him. Pilate saith unto them, Shall I crucify your King? The chief priests answered, We have no king but Cæsar. Then delivered he him therefore unto them to be crucified. And they took Jesus, and led him away. And he bearing his cross went forth into a place called the place of a skull, which is called in the Hebrew Golgotha: where they crucified him, and two other with him, on either side one, and Jesus in the midst. And Pilate wrote a title and put it on the cross. And the writing was, \divineName{Jesus of Nazareth the King of the Jews}. This title then read many of the Jews: for the place where Jesus was crucified was nigh to the city; and it was written in Hebrew, and Greek, and Latin. Then said the chief priests of the Jews to Pilate, Write not, The King of the Jews; but that he said, I am King of the Jews. Pilate answered, What I have written I have written. Then the soldiers, when they had crucified Jesus, took his garments, and made four parts, to every soldier a part; and also his coat: now the coat was without seam, woven from the top throughout. They said therefore among themselves, Let us not rend it, but cast lots for it, whose it shall be: that the scripture might be fulfilled, which saith, They parted my raiment among them, and for my vesture they did cast lots. These things therefore the soldiers did. Now there stood by the cross of Jesus his mother, add his mother's sister, Mary the wife of Cleophas, and Mary Magdalene. When Jesus therefore saw his mother, and the disciple standing by, whom he loved, he saith unto his mother, Woman, behold thy son! Then saith he to the disciple, Behold thy mother! And from that hour that disciple took her unto his own home. After this, Jesus knowing that all things were now accomplished, that the scripture might be fulfilled, saith, I thirst. Now there was set a vessel full of vinegar: and they filled a spunge with vinegar, and put it upon hyssop, and put it to his mouth. When Jesus therefore had received the vinegar, he said, It is finished: and he bowed his head, and gave up the ghost. \inrub{Here genuflect, and pause awhile.} The Jews therefore, because it was the preparation, that the bodies should not remain upon the cross on the sabbath day, (for that sabbath day was an high day,) besought Pilate that their legs might be broken, and that they might be taken away. Then came the soldiers, and brake the legs of the first, and of the other which was crucified with him. But when they came to Jesus, and saw that he was dead already, the brake not his legs: but one of the soldiers with a spear pierced his side, and forthwith came there out blood and water. And he that saw it bare record, and his record is true: and he knoweth that he saith true, that ye might believe. For these things were done, that the scripture should be fulfilled, A bone of him shall not be broken. And again another scripture saith, They shall look on him whom they pierced.

\subsubsection{Solemn Collect for the Catholic Church}

\begin{rubric}
	Then the Priest, standing at the Epistle side of the Altar, begins at once with joined hands:
\end{rubric}
\par\noindent

\gregorioscore{resources/gabc/Propers/GoodFriday/ChurchCollect.gabc}
%Let us pray dearly beloved, for the holy Church of God: that our God and Lord would vouchsafe to give her peace and unity, and preserve her throughout all the world: making subject unto her principalities and powers: and grant that, leading a quiet and peaceful life, we may glorify God the Father almighty.

\gregorioscore{resources/gabc/Propers/GoodFriday/Arise.gabc}

%\elcol{℣. Let us pray.}{℣. Orémus.}
%\elcol{℣. Let us bow the knee.}{℣. Flectámus génua.}
%\elcol{℣. Arise.}{℣. Leváte.}

%MANUAL ADJUSTMENT:
\clearpage
\begin{rubric}
	The Collect is sung in the ferial tone of the Collect of the Mass, with hands extended. And the same manner is observed in all that follow.
\end{rubric}

\lett{A}{lmighty} and everlasting God, who in Christ hast revealed thy glory to all nations: preserve the works of thy mercy; that thy Church, spread abroad over the whole world, may with steadfast faith persevere in the confession of thy name. Through the same.%\\

\subsubsection{Solemn Collect for the Chief Bishop}

\par\noindent
\gregorioscore{resources/gabc/Propers/GoodFriday/ChiefBishopCollect.gabc}
%Let us pray also for our most blessed Father \textit{N.}: that our God and Lord, who hath chosen him unto the order of bishops, may preserve him in health and safety to his holy Church, for the governance of the holy people of God.

\elcol{℣. Let us pray.}{℣. Orémus.}
\elcol{℣. Let us bow the knee.}{℣. Flectámus génua.}
\elcol{℣. Arise.}{℣. Leváte.}
%MANUAL ADJUSTMENT:
%
\lett{A}{lmighty} and everlasting God, by whose judgement all things are established: mercifully regard our prayers, and in thy goodness preserve him whom thou hast chosen to be our Bishop; that the Christian people who are governed by thine authority may under so great a Pontiff increase in the merits of their faith. Through.%\\

\subsubsection{Solemn Collect for All Ranks of Clergy}

\par\noindent
\gregorioscore{resources/gabc/Propers/GoodFriday/ClergyCollect.gabc}
%Let us pray also for all Bishops, Priests, and Deacons, for all Subdeacons, Acolytes, Exorcists, Door-keepers, Confessors, Readers, Virgins, and Widows: and for all the holy people of God.

\elcol{℣. Let us pray.}{℣. Orémus.}
\elcol{℣. Let us bow the knee.}{℣. Flectámus génua.}
\elcol{℣. Arise.}{℣. Leváte.}

\lett{A}{lmighty} and everlasting God, by whose Spirit the whole body of the Church is governed and sanctified: receive our supplications, which we offer before thee for all orders of the same; that by the bounty of thy grace they may faithfully serve thee in their several estates. Through . . . in the unity of the same.%\\

\subsubsection{Solemn Collect for the Roman Emperor}

\par\noindent
\gregorioscore{resources/gabc/Propers/GoodFriday/EmperorCollect.gabc}
%Let us pray also for our most Christian Emperor (elect) \textit{N.}: that our God and Lord would make subject unto him all the nations of the heathen, for our perpetual peace.

\elcol{℣. Let us pray.}{℣. Orémus.}
\elcol{℣. Let us bow the knee.}{℣. Flectámus génua.}
\elcol{℣. Arise.}{℣. Leváte.}
%
\lett{A}{lmighty} and everlasting God, in whose hand is the dominion of all things and the government of all kingdoms: graciously look upon the Roman Empire; that the nations which trust in their own fierceness may be overthrown by the right hand of thy power. Through.%\\

\subsubsection{Solemn Collect for the Catechumens}

\par\noindent
\gregorioscore{resources/gabc/Propers/GoodFriday/CatechumensCollect.gabc}
%Let us pray also for our catechumens: that our God and Lord would open the ears of their hearts, and the gate of mercy: that, receiving in the waters of regeneration the remission of all their sins, they also may be found in Christ Jesus our Lord.

\elcol{℣. Let us pray.}{℣. Orémus.}
\elcol{℣. Let us bow the knee.}{℣. Flectámus génua.}
\elcol{℣. Arise.}{℣. Leváte.}

\lett{A}{lmighty} and everlasting God, who dost continually enrich thy Church with a new offspring: increase the faith and understanding of our catechumens; that they, being born again in the water of baptism, may be numbered among the sons of thine adoption. Through.%\\

\subsubsection{Solemn Collect for Pilgrims}

\par\noindent
\gregorioscore{resources/gabc/Propers/GoodFriday/PilgrimCollect.gabc}
%Let us pray, dearly beloved, unto God the Father Almighty, that he would purge the world from all errors: would take away diseases: drive away famine: open the prisons: loosen the chains: grant unto pilgrims a safe return: to the sick, healing: and to them that travel by sea or air a haven of safety.

\elcol{℣. Let us pray.}{℣. Orémus.}
\elcol{℣. Let us bow the knee.}{℣. Flectámus génua.}
\elcol{℣. Arise.}{℣. Leváte.}

\lett{A}{lmighty} and everlasting God, the comfort of them that mourn, the strength of them that travail: let the prayers of them that cry out of any tribulation ascend unto thee; that in their necessities all may rejoice in the succour of thy loving kindness. Through.%\\

%MANUAL ADJUSTMENT:
\clearpage
\subsubsection{Solemn Collect for Heretics \& Schismatics}

\par\noindent
\gregorioscore{resources/gabc/Propers/GoodFriday/HereticksCollect.gabc}
%Let us pray also for heretics and schismatics: that our God and Lord would deliver them from all their errors; and vouchsafe to call them back to their holy mother, the Catholic and Apostolic Church.

\elcol{℣. Let us pray.}{℣. Orémus.}
\elcol{℣. Let us bow the knee.}{℣. Flectámus génua.}
\elcol{℣. Arise.}{℣. Leváte.}

\lett{A}{lmighty} and everlasting God, who savest all men, and wouldest not that any should perish: look upon the souls that are deceived by the craft of the devil; that the hearts of them that are gone astray being delivered from all perversity of heresy, may turn to wisdom and come again to the unity of thy truth. Through.%\\

\subsubsection{Solemn Collect for the Jews}

\par\noindent
\gregorioscore{resources/gabc/Propers/GoodFriday/JewsCollect.gabc}
%Let us pray also for the faithless Jews: that our God and Lord would take away the veil from their hearts; that they also may acknowledge Jesus Christ, our Lord.

\begin{rubric}
	\blackrubric{Amen} is not said, nor \blackrubric{Let us pray,} nor \blackrubric{Let us bow the knee,} nor \blackrubric{Arise,} but at once is said:
\end{rubric}

\lett{A}{lmighty} and everlasting God, who deniest not thy mercy even to the faithless Jews: graciously hear our prayers, which we offer for the blindness of this people; that they acknowledging the light of thy truth, which is Christ, may be delivered from their darkness. Through the same.%\\

\subsubsection{Solemn Collect for Infidels}

\par\noindent
\gregorioscore{resources/gabc/Propers/GoodFriday/HeathenCollect.gabc}
%Let us pray also for the heathen: that God Almighty would take away the iniquity from their hearts; that, forsaking their idols, they may be turned unto the living and true God, and to his only Son, Jesus Christ, our God and Lord.

\elcol{℣. Let us pray.}{℣. Orémus.}
\elcol{℣. Let us bow the knee.}{℣. Flectámus génua.}
\elcol{℣. Arise.}{℣. Leváte.}

\lett{A}{lmighty} and everlasting God, who desirest not the death of sinners but rather that they should live: mercifully receive our prayer, and deliver them from the worship of idols; and gather them unto thy holy Church, to the praise and glory of thy name. Through.

\subsubsection{Veneration of the Cross}
\begin{rubric}
	Having finished the Prayers, the Priest lays aside his Chasuble and goes to the Epistle side of the Altar, where at the back corner he receives from the Deacon the Cross which has previously been made ready on the Altar. With his face to the people, he uncovers the upper portion, beginning alone \blackrubric{Behold the wood of the Cross,} and from thence he is helped in singing the rest by the Ministers until \blackrubric{O come, let us worship.}\par
When the Choir sing \blackrubric{O come, let us worship,} all prostrate themselves except the Celebrant. He then advances to the front corner of the same Epistle side: and uncovering the right arm of the Cross, and raising it a little, he begins in a higher key than at first: \blackrubric{Behold the wood of the Cross,} the others singing and worshipping as before. Then the Priest proceeds to the middle of the Altar: and uncovering the Cross entirely and raising it, for the third time he begins in a still higher key: \blackrubric{Behold the wood of the Cross,} the others singing and worshipping as before.
\end{rubric}

\gregorioscore{resources/gabc/Propers/GoodFriday/Behold.gabc}
%Behold the wood of the Cross, whereon was hung the world's salvation.

%\textit{Choir.} ℟. O come, let us worship.

\begin{rubric}
	Then the Priest alone carries the Cross to the place prepared before the Altar, and kneeling there he puts it down: then laying aside his shoes, he proceeds to worship the Cross %ADDITION TO CLARIFY MANNER OF WORSHIP:
	with the worship of 
	\begin{otherlanguage}{greek}
	προσκύνησις, 
	\end{otherlanguage}
	that is, dulia.
	\par
	And kneeling thrice before, he kisses it. This done, he returns, and puts on his shoes and the Chasuble. Afterwards the Ministers of the Altar, then the other Clerics and laics, two and two, kneeling thrice, as has been said, worship the Cross. Meantime, while the worship of the Cross is proceeding, the Reproaches are sung, together with the other things that follow, either all or in part, according as the many or few that worship the Cross require: which also the Priest sitting at the bench reads with the Ministers, as below:
\end{rubric}

\textsc{Celebrant.} O my people, what have I done unto thee? or wherein have I wearied thee? Answer me.\par
℣. Because I brought thee forth from the land of Egypt: thou hast prepared a Cross for thy Saviour.

\textsc{Deacon.} \begin{otherlanguage}{greek} Ἅγιος ὁ Θεός.\end{otherlanguage} (Hagios ho Theós.)\par
\textsc{Subdeacon.} Holy God.\par
\textsc{Deacon.} \begin{otherlanguage}{greek} Ἅγιος ἰσχυρός.\end{otherlanguage} (Hagios ischyrós.)\par
\textsc{Subdeacon.} Holy Mighty.\par
\textsc{Deacon.} \begin{otherlanguage}{greek} Ἅγιος ἀθάνατος, ἐλέησον ἡμᾶς.\end{otherlanguage} (Hagios athánatos, eléison himas.)\par
\textsc{Subdeacon.} Holy Immortal, have mercy upon us.\par
\textsc{Celebrant.} Because I led thee through the desert forty years, and fed thee with manna, and brought thee into a land exceeding good: thou hast prepared a Cross for thy Saviour.\par
\textsc{Deacon.} \begin{otherlanguage}{greek} Ἅγιος ὁ Θεός.\end{otherlanguage} (Hagios ho Theós.)\par
\textsc{Subdeacon.} Holy God.\par
\textsc{Deacon.} \begin{otherlanguage}{greek} Ἅγιος ἰσχυρός.\end{otherlanguage} (Hagios ischyrós.)\par
\textsc{Subdeacon.} Holy Mighty.\par
\textsc{Deacon.} \begin{otherlanguage}{greek} Ἅγιος ἀθάνατος, ἐλέησον ἡμᾶς.\end{otherlanguage} (Hagios athánatos, eléison himas.)\par
\textsc{Subdeacon.} Holy Immortal, have mercy upon us.\par
\textsc{Celebrant.} What more could I have done for thee that I have not done? I indeed did plant thee, my vineyard, exceeding fair: and thou art become very bitter unto me: for vinegar thou gavest to quench my thirst: and hast pierced with a spear the side of thy Saviour.\par
\textsc{Deacon.} \begin{otherlanguage}{greek} Ἅγιος ὁ Θεός.\end{otherlanguage} (Hagios ho Theós.)\par
\textsc{Subdeacon.} Holy God.\par
\textsc{Deacon.} \begin{otherlanguage}{greek} Ἅγιος ἰσχυρός.\end{otherlanguage} (Hagios ischyrós.)\par
\textsc{Subdeacon.} Holy Mighty.\par
\textsc{Deacon.} \begin{otherlanguage}{greek} Ἅγιος ἀθάνατος, ἐλέησον ἡμᾶς.\end{otherlanguage} (Hagios athánatos, eléison himas.)\par
\textsc{Subdeacon.} Holy Immortal, have mercy upon us.\par
\textsc{Celebrant.} I did scourge Egypt with her first-born for thy sake: and thou hast scourged me and delivered me up.\par
\textsc{Deacon and Subdeacon.} O my people, what have I done unto thee, or wherein have I wearied thee? Answer me.\par 
\textsc{Celebrant.} I led thee out of Egypt, drowning Pharaoh in the Red Sea: and thou hast delivered me unto the chief priests.\par
\textsc{Deacon and Subdeacon.} O my people, what have I done unto thee, or wherein have I wearied thee? Answer me.\par
\textsc{Celebrant.} I opened the sea before thee: and thou hast opened my side with a spear.\par
\textsc{Deacon and Subdeacon.} O my people, what have I done unto thee, or wherein have I wearied thee? Answer me.\par
\textsc{Celebrant.} I went before thee in a pillar of cloud: and thou hast led me unto the judgment hall of Pilate.\par
\textsc{Deacon and Subdeacon.} O my people, what have I done unto thee, or wherein have I wearied thee? Answer me.\par
\textsc{Celebrant.} I fed thee with manna in the desert: and thou hast stricken me with blows and scourges.\par
\textsc{Deacon and Subdeacon.} O my people, what have I done unto thee, or wherein have I wearied thee? Answer me.\par
\textsc{Celebrant.} I gave thee to drink of the water of salvation from the rock: and thou hast given me gall and vinegar to drink.\par
\textsc{Deacon and Subdeacon.} O my people, what have I done unto thee, or wherein have I wearied thee? Answer me.\par
\textsc{Celebrant.} For thee I smote the kings of the Canaan-ites: and thou hast smitten my head with a reed.\par
\textsc{Deacon and Subdeacon.} O my people, what have I done unto thee, or wherein have I wearied thee? Answer me.\par
\textsc{Celebrant.} I gave thee a royal sceptre: and thou hast given unto my head a crown of thorns.\par
\textsc{Deacon and Subdeacon.} O my people, what have I done unto thee, or wherein have I wearied thee? Answer me.\par
\textsc{Celebrant.} I exalted thee with great power: and thou hast hanged me upon the gibbet of the Cross.\par
\textsc{Deacon and Subdeacon.} O my people, what have I done unto thee, or wherein have I wearied thee? Answer me.\par
\textsc{Celebrant.} We worship thy Cross, O Lord: and praise and glorify thy holy resurrection: for behold, by virtue of the tree joy hath come to the whole world. \textit{Ps.} God be merciful unto us and bless us :\par
\textsc{Deacon and Subdeacon.} And shew us the light of his countenance, and be merciful unto us.\par
\textsc{Celebrant.} We worship thy Cross, O Lord: and praise and glorify thy holy resurrection: for behold, by virtue of the tree joy hath come to the whole world.\par
\textsc{Deacon and Subdeacon.} Faithful Cross, above all other, one and only noble tree: none in foliage, none in blossom, none in fruit thy peer may be.\par
Sweetest wood and sweetest iron, sweetest weight is hung on thee.\par
\textsc{Celebrant.} \textit{Hymn.} Sing, my tongue, the glorious battle, sing the ending of the fray, now above the Cross, loud triumphant lay: tell how Christ, the world's Redeemer, as a victim won the day.\par
\textsc{Deacon and Subdeacon.} Faithful Cross, above all other, one and only noble tree: none in foliage, none in blossom, none in fruit thy peer may be.\par
\textsc{Celebrant.} God in pity saw man fallen, shamed and sunk in misery, when he fell on death by tasting fruit of the forbidden tree: then another tree was chosen which the world from death should free.\par
\textsc{Deacon and Subdeacon.} Sweetest wood and sweetest iron, sweetest weight is hung on thee.\par
\textsc{Celebrant.} Thus the scheme of our salvation, was of old in order laid: that the manifold deceiver's art by art might be outweighed: and the lure the toe put forward into means of healing made.\par
\textsc{Deacon and Subdeacon.} Faithful Cross, above all other, one and only noble tree: none in foliage, none in blossom, none in fruit thy peer may be.\par
\textsc{Celebrant.} Therefore when the appointed fulness of the holy time was come, he was sent who maketh all things forth from God's eternal home; thus he came to earth, incarnate, offspring of a virgin's womb.\par
\textsc{Deacon and Subdeacon.} Sweetest wood and sweetest iron, sweetest weight is hung on thee.\par
\textsc{Celebrant.} Lo! he lies, an Infant weeping, where the narrow manger stands, while the Mother-Maid his members wraps in mean and lowly bands, and the swaddling clothes is winding round his helpless feet and hands.\par
\textsc{Deacon and Subdeacon.} Faithful Cross, above all other, one and only noble tree: none in foliage, none in blossom, none in fruit thy peer may be.\par
\textsc{Celebrant.} Thirty years among us dwelling, his appointed time fulfilled, born for this, he meets his Passion, for that this he freely willed, on the Cross the Lamb is lifted where his life-blood shall be spilled.\par
\textsc{Deacon and Subdeacon.} Sweetest wood and sweetest iron, sweetest weight is hung on thee.\par
\textsc{Celebrant.} He endured the nails, the spitting, vinegar, and spear, and reed; from that holy body broken blood and water forth proceed: earth, and stars, and sky, and ocean by that flood from stain are freed.\par
\textsc{Deacon and Subdeacon.} Faithful Cross, above all other, one and only noble tree: none in foliage, none in blossom, none in fruit thy peer may be.\par
\textsc{Celebrant.} Bend thy boughs, O tree of glory! thy relaxing sinews bend; for awhile the ancient rigour that thy birth bestowed, suspend; and the King of heavenly beauty on thy bosom gently tend!\par
\textsc{Deacon and Subdeacon.} Sweetest wood and sweetest iron, sweetest weight is hung on thee.\par
\textsc{Celebrant.} Thou alone wast counted worthy this world's ransom to uphold ; for a shipwreck'd race preparing harbour, like the ark of old; with the sacred blood anointed from the smitten Lamb that rolled.\par
\textsc{Deacon and Subdeacon.} Faithful Cross, above all other, one and only noble tree: none in foliage, none in blossom, none in fruit thy peer may be.\par
\textsc{Celebrant.} To the Trinity be glory everlasting, as is meet; equal to the Father, equal to the Son, and Paraclete: Trinal Unity, whose praises all created things repeat. Amen.\par
\textsc{Deacon and Subdeacon.} Sweetest wood and sweetest iron, sweetest weight is hung on thee.\par

\begin{rubric}
	Towards the end of the worship of the Cross candles are lighted upon the Altar: and the Deacon, receiving the Burse with the Corporal, spreads the Corporal in the usual way, and puts a Purificator by it; and when the adoration is over, reverently receives the Cross and carries it back to the Altar. Then a Procession is formed to the place where on the previous day the Sacrament bad been laid. A Subdeacon goes first with the Cross between two Acolytes bearing candlesticks with lighted candles, and then the Clergy in order, last of all the Priest with the Ministers. When they arrive at the place of the Sacrament, torches are kindled, which are not extinguished till after the consumption of the Sacrament. The Priest genuflects before the Sacrament and prays awhile: the Deacon meanwhile opens the receptacle in which lies hidden the Body of the Lord. Then the Priest, rising, puts incense into two thuribles without blessing, the Deacon presenting the boat, and then kneeling he censes the Sacrament.\par
	The Deacon takes the Chalice with the Sacrament from the receptacle and gives it into the hands of the Priest, covering it with the ends of the Veil which is around his shoulders. Then they proceed in the order in which they came: a baldachino is carried over the Sacrament: and two Acolytes with thuribles continually cense the Sacrament. Meanwhile is sung the following hymn:
\end{rubric}

\subsubsection{Vexilla Regis prodeunt}

\gregorioscore{resources/gabc/Propers/GoodFriday/RoyalBanners.gabc}

\subsubsection{Liturgy of the Presanctified Gifts}
\begin{rubric}
	When the Priest arrives at the Altar, he places the Chalice upon it, genuflects, and again censes it, and going up takes the Host out of the Chalice and places it upon the Paten which the Deacon holds: and receiving the Paten from the hand of the Deacon, he places the sacred Host upon the Corporal, saying nothing. If he touches the Sacrament, he washes his hands in some vessel. Meanwhile the Deacon puts wine in the Chalice, and the Subdeacon water, which the Priest does not bless, nor does he say over it the usual Prayer: but receiving the Chalice from the Deacon, be places it upon the Altar, saying nothing: and the Deacon covers it with the Pall.\par
Then he puts incense in the thurible without blessing it, and censes the Oblations, Cross, and Altar in the accustomed manner, genuflecting before and after, and whenever be crosses before the Sacrament.
\end{rubric}

\begin{rubric}
	When he censes the Altar, he says:
\end{rubric}
\lett{L}{et} my prayer, O Lord, be set forth in thy sight as the incense: and let the lifting up of my hands be an evening sacrifice. Set a watch, O Lord, before my mouth, and keep the door of my lips: O let not mine heart be inclined to any evil thing, let me not be occupied in ungodly works.

\begin{rubric}
	When he gives back the thurible to the Deacon he says:
\end{rubric}

\lett{T}{he} Lord kindle in us the fire of his love and the flame of everlasting charity. Amen.

\lettspace

\begin{rubric}
	And he himself is not censed.
\end{rubric}

\begin{rubric}
	Afterwards he washes his hands a little beyond the Altar on the Epistle side, saying nothing: then bowing in the middle of the Altar, with hands joined, he says:
\end{rubric}

\lett{I}{n} a humble spirit and with a contrite heart may we be accepted of thee, O Lord: and so let our sacrifice be offered in thy sight this day, that it may be pleasing unto thee, O Lord God.

%MANUAL ADJUSTMENT:
\clearpage
\begin{rubric}
	Then turning to the people at the Gospel corner, he says as usual:
\end{rubric}

\lett{P}{ray,} brethren, that my sacrifice and yours may be acceptable to God the Father almighty.

\lettspace

\begin{rubric}
	And he turns back by the same way, not completing the circle: and, omitting the rest, proceeds:
\end{rubric}
\par\noindent
Let us pray: Commanded by saving precepts, and taught by divine institution, we are bold to say:

\lett{O}{ur} Father, who art in heaven: Hallowed be thy name: Thy kingdom come: Thy will be done, on earth as it is in heaven. Give us this day our daily bread: And forgive us our trespasses, as we forgive those that trespass against us. And lead us not into temptation.\par
℟. But deliver us from evil.

\begin{rubric}
	The Priest says silently \blackrubric{Amen.} In the same voice in which he said \blackrubric{Our Father,} and without \blackrubric{Let us pray,} he says in the tone of the Collect of the Ferial Mass:
\end{rubric}

\lett{D}{eliver} us, O Lord, we beseech thee, from all evils, past, present, and to come: and at the intercession of the blessed and glorious ever Virgin Mary Mother of God, with thy blessed Apostles Peter and Paul, and with Andrew, and all the Saints, \inrub{he does not sign himself with the Paten} favourably grant peace in our days: that by the help of thine availing mercy we may ever both be free from sin, and safe from all distress. Through the same thy Son Jesus Christ our Lord: Who liveth and reigneth with thee, in the unity of the Holy Ghost, God, throughout all ages, world without end.\par
℟. Amen.

%\subsubsection{Confession}
\begin{rubric}
    Then the Deacon turneth to the People and saith to those who come to receive the Holy Communion,
\end{rubric}
\lett{Y}{e} who do truly and earnestly repent you of your sins, and are in love and charity with your neighbours, and intend to lead a new life, following the commandments of God, and walking from henceforth in his holy ways; Draw near with faith, and take this holy Sacrament to your comfort; and make your humble confession to Almighty God, devoutly kneeling.
\begin{rubric}
Then shall this General Confession be made, in an audible voice, by the Priest, bowing, and all those who are minded to receive the Holy Communion, humbly kneeling.
\end{rubric}

\lett{A}{lmighty} God, Father of our Lord Jesus Christ, Maker of all things, Judge of all men; We acknowledge and bewail our manifold sins and wickedness, Which we, from time to time, most grievously have committed, By thought, word, and deed, Against thy Divine Majesty, Provoking most justly thy wrath and indignation against us. We do earnestly repent, And are heartily sorry for these our misdoings; The remembrance of them is grievous unto us; The burden of them is intolerable. Have mercy upon us, Have mercy upon us, most merciful Father; For thy Son our Lord Jesus Christ's sake, Forgive us all that is past; And grant that we may ever hereafter Serve and please thee in newness of life, To the honour and glory of thy Name; Through Jesus Christ our Lord. Amen.
\begin{rubric}
	Then shall the Priest stand up and, turning to the People, say,
\end{rubric}

\lett{A}{lmighty} God, our heavenly Father, who of his great mercy hath promised forgiveness of sins to all those who with hearty repentance and true faith turn unto him; Have mercy upon you; pardon {\ding{64}} and deliver you from all your sins; confirm and strengthen you in all goodness; and bring you to everlasting life; through Jesus Christ our Lord. \textit{Amen.}

%\subsubsection{Comfortable Words}
\begin{rubric}
Then shall the Priest, facing the People, say,
\end{rubric}\noindent
\begin{center}
	\textsc{Hear what comfortable words our Saviour Christ saith unto all who truly turn to him.}
\end{center}
\par\noindent
Come unto me, all ye that travail and are heavy laden, and I will refresh you.
\par\noindent
    So God loved the world, that he gave his only-begotten Son, to the end that all that believe in him should not perish, but have everlasting life.
    \par\noindent
    \begin{center}
		\textsc{Hear also what Saint Paul saith.}
	\end{center}
\par\noindent
    This is a true saying, and worthy of all men to be received, That Christ Jesus came into the world to save sinners. 

    \begin{center}
		\textsc{Hear also what Saint John saith.}
	\end{center}
    \par\noindent
    If any man sin, we have an Advocate with the Father, Jesus Christ the righteous; and he is the Propitiation for our sins.

%\subsubsection{Prayer of Humble Access}
\begin{rubric}
    The Priest alone then saith, in an audible voice,
\end{rubric}
\lett{W}{e} do not presume to come to this thy Table, O merciful Lord, trusting in our own righteousness, but in thy manifold and great mercies. We are not worthy so much as to gather up the crumbs under thy Table. But thou art the same Lord, whose property is always to have mercy: Grant us therefore, gracious Lord, so to eat the flesh of thy dear Son Jesus Christ, and to drink his blood, that our sinful bodies may be made clean by his Body, and our souls washed through his most precious Blood, and that we may evermore dwell in him, and he in us.\par
℟. Amen.

%\subsubsection{Communion of the Priest}

\begin{rubric}
	Then the Celebrant, having made a reverence down to the ground, places the Paten beneath the Sacrament, which he takes in his right hand and elevates so that it can be seen by the people: and immediately he divides it over the Chalice into three parts, the last of which he puts into the Chalice in the usual way, saying nothing. \blackrubric{The Peace of the Lord} is not said, nor \blackrubric{O Lamb of God,} nor is the kiss of peace given. Then, omitting the first two Prayers, he says only the following:
\end{rubric}

\lett{L}{et} the partaking of thy Body, O Lord Jesu Christ, which I, unworthy, presume to receive, turn not to my judgement and condemnation: but of thy goodness let it avail unto me for protection of mind and of body, that I may receive thy healing: Who livest and reignest with God the Father in the unity of the Holy Ghost, God, throughout all ages, world without end. Amen.

\begin{rubric}
	Then he genuflects and receives the Paten with the Body of Christ: and with the deepest humility and reverence says:
\end{rubric}

I will receive the bread of heaven, and call upon the name of the Lord.

\begin{rubric}
	He beats his breast, saying thrice:
\end{rubric}

Lord, I am not worthy that thou shouldest come under my roof: but speak the word only, and my soul shall be healed.

\begin{rubric}
	Then he signs himself with the Sacrament, saying:
\end{rubric}

The Body of our Lord Jesus Christ preserve my soul unto everlasting life. Amen.

\begin{rubric}
	And he reverently consumes the Body.
\end{rubric}

\begin{rubric}
	Then, omitting everything which is wont to be said before the consumption of the Blood, he at once reverently consumes the particle of the Host with the Chalice.
\end{rubric}

%\subsubsection{Communion of the People}
\begin{rubric}
	The Priest then, turning towards the Congregation, holds a particle over the Chalice and says:
\end{rubric}
\elcol{℣. Behold the Lamb of God. Behold him who taketh away the sins of the world.}{℣. Ecce Agnus Dei, ecce, qui tollit peccáta mundi.}
\elcol{℟. Lord, I am not worthy that thou shouldest come under my roof, but speak the word only, and my soul shall be healed.}{℟. Dómine, non sum dignus, ut intres sub tectum meum, sed tantum dic verbo, et sanábitur ánima mea.}
\elcol{℟. Lord, I am not worthy that thou shouldest come under my roof, but speak the word only, and my soul shall be healed.}{℟. Dómine, non sum dignus, ut intres sub tectum meum, sed tantum dic verbo, et sanábitur ánima mea.}
\elcol{℟. Lord, I am not worthy that thou shouldest come under my roof, but speak the word only, and my soul shall be healed.}{℟. Dómine, non sum dignus, ut intres sub tectum meum, sed tantum dic verbo, et sanábitur ánima mea.}

\begin{rubric}
	The Priest then distributes Communion, in both kinds, to the People in the usual manner.
\end{rubric}

\begin{rubric}
	After distributing the Communion to the People, he makes the ablution of his fingers in the usual manner and taken the purification. He then says, bowing in the midst of the Altar, with hands joined:
\end{rubric}

\lett{G}{rant,} O Lord, that what we have taken with our mouths we may receive with a pure heart: and from a temporal gift may it become to us an everlasting remedy.

\begin{rubric}
	\blackrubric{May thy Body} is not said, nor the Postcommunion, nor \blackrubric{Let this my bounden duty,} nor is the blessing given: but having made a reverence to the Altar, the Priest retires with the Ministers.
\end{rubric}

%\subsubsection{Thanksgiving}
\begin{rubric}
    Standing with his hands joined, the Priest shall alone say the Thanksgiving, in an audible voice.
\end{rubric}
\letuspray
\lett{A}{lmighty} and everliving God, we most heartily thank thee, for that thou dost vouchsafe to feed us who have duly received these holy mysteries with the spiritual food of the most precious Body and Blood of thy Son our Saviour Jesus Christ; and dost assure us thereby of thy favour and goodness towards us; and that we are very members incorporate in the mystical body of thy Son, which is the blessed company of all faithful people; and are also heirs through hope of thy everlasting kingdom, by the merits of his most precious death and passion. And we humbly beseech thee, O heavenly Father, so to assist us with thy grace, that we may continue in that holy fellowship, and do all such good works as thou hast prepared for us to walk in; through Jesus Christ our Lord, to whom, with thee and the Holy Ghost, be all honour and glory, world without end.\par
℟. Amen.

\begin{rubric}
	Evensong is then said without singing: and the Altar is stripped.
\end{rubric}


